% Computer Science A --- HTTP (JSON), REST, and Javalin (\input inside \texttt{main.tex} after \texttt{cs-a/06-practice.tex}).
% HTTP free-response prompts at the front are migrated from \texttt{cs-a/03-practice.tex}.

\runningheader{Computer Science A}{RESTful APIs (Javalin)}{Page \thepage\ of \numpages}

\begingradingrange{csarest}

\uplevel{\par\textbf{HTTP}\par}

\question[4]
What is HTTP? Why is it important to know about?
\begin{solutionorbox}[1.25in]
  Hypertext Transfer Protocol: request/response protocol underlying the Web and many APIs. Important for building clients/servers, debugging APIs, caching, auth, and status handling.
\end{solutionorbox}

\question[4]
What are common HTTP verbs used when a client application is making a request?
\begin{solutionorbox}[1.15in]
  \texttt{GET}, \texttt{POST}, \texttt{PUT}, \texttt{PATCH}, \texttt{DELETE}, often \texttt{HEAD}/\texttt{OPTIONS}; semantics express safe/idempotent reads vs resource creation/update/removal.
\end{solutionorbox}

\question[4]
What are some common HTTP status codes that can be included in a response?
\begin{solutionorbox}[1.35in]
  Examples: \texttt{200} OK, \texttt{201} Created, \texttt{204} No Content, \texttt{400} Bad Request, \texttt{401}/\texttt{403} auth/forbidden, \texttt{404} Not Found, \texttt{500} server error---clients branch on these.
\end{solutionorbox}

\newpage

\uplevel{\par\textbf{HTTP (Hypertext Transfer Protocol)}\par}


\question[4]
What does \textbf{API} stand for?
\begin{choices}
  \choice Advanced Programming Interface
  \choice Application Programming Infrastructure
  \CorrectChoice Application Programming Interface
  \choice Apply Programming Interfaces
\end{choices}
\begin{solution}
  \textbf{API} expands to \textbf{Application Programming Interface}: the documented surface (endpoints, types, auth) through which programs call another system.
\end{solution}

\question[4]
What is JSON? When would we work with JSON in our application?
\begin{solutionorbox}[1.00in]
  JavaScript Object Notation: text format for nested objects/arrays. Common for REST request/response bodies, config, and browser--server data exchange.
\end{solutionorbox}

\question[4]
What is XML? When would we work with XML in our application?
\begin{solutionorbox}[1.00in]
  Extensible Markup Language: text format with nested elements and attributes (often a single root). Used for some REST/SOAP payloads, configuration, feeds, and interoperability with systems that already standardize on XML.
\end{solutionorbox}


\question[4]
What is HTML? When would we work with HTML in our application?
\begin{solutionorbox}[1.00in]
  HyperText Markup Language: markup for documents with structure, links, and embedded media, meant for browsers to render. Typical for web pages and some hypermedia-style APIs; less often the main machine-to-machine API body than JSON, but common for human-readable responses or mixed content.
\end{solutionorbox}


\newpage

\question[4]
Which of these is \textbf{not} a common \uline{state representation format} for data sent to or returned from a REST-style HTTP API?
\begin{choices}
  \choice JSON
  \choice XML
  \choice HTML
  \CorrectChoice CSS
\end{choices}
\begin{solution}
  JSON, XML, and HTML can carry resource representations over HTTP. 
  \uline{CSS} styles documents; 
  it is not used as the primary API payload format for resource state.
\end{solution}

\question[4]
HTTP responses with status codes in the \textbf{100--199} range are generally classified as which category?
\begin{choices}
  \CorrectChoice Informational
  \choice Success
  \choice Redirection
  \choice Client error
\end{choices}
\begin{solution}
  The 1xx class is informational (e.g.\ request received, processing continues).
\end{solution}

\question[4]
HTTP responses with status codes in the \textbf{200--299} range are generally classified as which category?
\begin{choices}
  \choice Informational
  \CorrectChoice Success
  \choice Redirection
  \choice Client error
\end{choices}
\begin{solution}
  The 2xx class indicates successful processing (\texttt{200 OK}, \texttt{201 Created}, etc.).
\end{solution}


\question[4]
HTTP responses with status codes in the \textbf{300--399} range are generally classified as which category?
\begin{choices}
  \choice Informational
  \choice Success
  \CorrectChoice Redirection
  \choice Client error
\end{choices}
\begin{solution}
  The 3xx class signals redirection or further action needed to complete the request.
\end{solution}



\question[4]
HTTP responses with status codes in the \textbf{400--499} range are generally classified as which category?
\begin{choices}
  \choice Informational
  \choice Success
  \choice Redirection
  \CorrectChoice Client error
\end{choices}
\begin{solution}
  The 4xx class indicates a problem with the request from the client's perspective (\texttt{404 Not Found}, \texttt{400 Bad Request}, etc.).
\end{solution}


\newpage

\question[4]
HTTP responses with status codes in the \textbf{500--599} range are generally classified as which category?
\begin{choices}
  \choice Informational
  \choice Success
  \choice Redirection
  \CorrectChoice Server error
\end{choices}
\begin{solution}
  The 5xx class indicates the server failed to fulfill an apparently valid request (\texttt{500}, \texttt{503}, etc.).
\end{solution}

\question[4]
Which option correctly describes two common HTTP methods and their typical roles?
\begin{choices}
  \choice GET submits data, and POST retrieves data
  \choice PUT retrieves data, and DELETE submits data
  \CorrectChoice PUT updates (or replaces) a resource, and DELETE removes a resource
  \choice None of the above
\end{choices}
\begin{solution}
  PUT commonly updates or replaces representations; DELETE removes resources. GET retrieves; POST often creates or submits data.
\end{solution}


\question[4]
Which HTTP method is commonly used for submitting HTML form data to the server?
\begin{choices}
  \choice GET
  \CorrectChoice POST
  \choice PUT
  \choice DELETE
\end{choices}
\begin{solution}
  POST carries a body and is the usual choice for form submissions (GET may be used for idempotent lookups with query strings).
\end{solution}



\question[4]
Which HTTP status code most directly indicates that the server is \emph{currently} unable to handle the request (e.g.\ overload or maintenance)?
\begin{choices}
  \choice 500
  \choice 404
  \CorrectChoice 503
  \choice 302
\end{choices}
\begin{solution}
  \texttt{503 Service Unavailable} signals temporary inability to handle traffic; \texttt{500} is a generic server error; \texttt{404} is client-visible ``not found.''
\end{solution}


\newpage

\question[4]
In the raw request below, what is the client primarily doing?
\begin{lstlisting}[style=http]
POST /echo/post/json HTTP/1.1
Host: oracle.com
Accept: application/json
Content-Type: application/json
Content-Length: 81
{
    "Id": 78912,
    "Customer": "Frederick Smith",
    "Quantity": 1,
    "Price": 20.00
}
\end{lstlisting}
\begin{choices}
  \choice Retrieving a JSON object from the server
  \CorrectChoice Sending a body to create or add a resource (or invoke a processing endpoint)
  \choice Updating an existing resource with PUT semantics
  \choice None of the above
\end{choices}
\begin{solution}
  A \texttt{POST} with a JSON body typically submits data for the server to accept or create; it is not a retrieval.
\end{solution}



\question[4]
What does this HTTP response indicate?
\begin{lstlisting}[style=http]
HTTP/1.1 404 Not Found
Date: Sun, 18 Oct 2012 10:36:20 GMT
Server: Apache/2.2.14 (Win32)
Content-Length: 230
Connection: Closed
Content-Type: text/html; charset=iso-8859-1
<!DOCTYPE HTML PUBLIC "-//IETF//DTD HTML 2.0//EN">
<html>
<head>
    <title>404 Not Found</title>
</head>
<body>
    <h1>Not Found</h1>
    <p>The requested URL /t.html was not found on this server.</p>
</body>
</html>
\end{lstlisting}
\begin{choices}
  \CorrectChoice The requested resource was not found (client-visible error; 4xx)
  \choice There was a server-side error (5xx)
  \choice The response returned the expected successful payload
  \choice The server redirected the client to another URL (3xx)
\end{choices}
\begin{solution}
  \texttt{404 Not Found} is a client error: the URI did not map to a resource on the server.
\end{solution}


\newpage

\question[4]
What does this HTTP response indicate?
\begin{lstlisting}[style=http]
HTTP/1.1 200 OK
Access-Control-Allow-Origin: *
Content-Type: text/html
<html>
    <body>
        <h1>Hello World</h1>
    <body>
</html>
\end{lstlisting}
\begin{choices}
  \choice There was an error on the server side
  \choice The server redirected the client to another URL
  \choice The request was invalid and could not be completed
  \CorrectChoice The request succeeded and the body is an HTML document
\end{choices}
\begin{solution}
  \texttt{200 OK} with \texttt{text/html} means success; the payload is HTML for the client to render.
\end{solution}

\question[4]
An HTTP \emph{request} message does not include a status line with a numeric status code (that appears on the response).
\begin{choices}
  \CorrectChoice True
  \choice False
\end{choices}
\begin{solution}
  Requests have a request line (method, path, version); status codes belong in the response status line.
\end{solution}


\question[4]
What is correct about HTTP status codes?
\begin{choices}
  \choice They range from the 100s through the 900s as standard classes
  \CorrectChoice They tell the client whether the exchange succeeded, failed, or needs follow-up (and often why)
  \choice They are unimportant for API development
  \choice They are defined by CSS
\end{choices}
\begin{solution}
  Standard classes are 1xx--5xx; they are central to debugging clients, servers, and proxies.
\end{solution}

\question[4]
If a response has status \texttt{500 Internal Server Error}, that typically means an unexpected error occurred on the server; such outcomes should be rare and investigated in production.
\begin{choices}
  \CorrectChoice True
  \choice False
\end{choices}
\begin{solution}
  \texttt{500} indicates a server failure; good operations and tests aim to avoid surfacing unhandled errors to clients.
\end{solution}


\newpage

\question[4]
In the code snippet below, what kind of payload is shown?
\begin{lstlisting}[style=http]
{
  "employee": {
    "name": "Alfred",
    "salary": 45000,
    "married": false
  }
}
\end{lstlisting}
\begin{choices}
  \CorrectChoice A JSON object
  \choice An XML document
  \choice An HTML page
  \choice A C\# object
\end{choices}
\begin{solution}
  The payload is JSON (object with nested \texttt{employee} object).
\end{solution}



\question[4]
Where is the inconsistency in the JSON structure below?
\begin{lstlisting}[style=http]
{
  "employees": [
    {
      "id": 1,
      "first_name": "Sebastian",
      "last_name": "Eschweiler",
      "email": "sebastianE@gmail.com"
    },
    {
      "id": 2,
      "first_name": "Steve",
      "lastName": "Palmer",
      "email": "steveP@gmail.com"
    }
  ]
}
\end{lstlisting}
\begin{choices}
  \choice in a value
  \CorrectChoice in a property
  \choice in an object
  \choice in a method
\end{choices}
\begin{solution}
  Property naming is inconsistent (\texttt{last\_name} vs.\ \texttt{lastName})---a schema/style issue among sibling objects.
\end{solution}


















\newpage


\uplevel{\par\textbf{REST}\par}

\question[4]
What is REST and what are its key constraints?
\begin{solutionorbox}[1.00in]
  \textbf{REST} (Representational State Transfer) is an \uline{architectural style} for building networked systems on top of HTTP: clients and servers exchange representations of \textbf{resources} through a small, uniform set of operations.

  \noindent\textbf{Typical constraints (as usually taught):}
  \begin{itemize}
    \item \textbf{Resources and URIs:} important things are named with stable identifiers. A \textbf{URI} (Uniform Resource Identifier, RFC~3986) is the general term; a \textbf{URL} is a URI that also \emph{locates} a resource (e.g.\ an \texttt{https://\ldots} address is both). HTTP targets a resource via a URI (path and query on the wire, often embedded in a full URL).
    \item \textbf{Stateless requests:} each call should carry what the server needs to understand it; the server does not rely on hidden per-client session state between requests.
    \item \textbf{Uniform interface:} use standard HTTP methods and status codes with consistent, documented semantics (\texttt{GET} read, \texttt{POST} create/submit, \texttt{PUT}/\texttt{PATCH} update, \texttt{DELETE} remove, etc.).
    \item \textbf{Representations:} payloads are serializations of a resource (often JSON or XML), not the underlying storage row-for-row.
  \end{itemize}

  \noindent\textbf{Note:} Production APIs often relax or extend pure REST; the goals are interoperability, predictability, and clear separation of concerns.
\end{solutionorbox}

\question[4]
\textbf{True or false:} REST (\textbf{RE}presentational \textbf{S}tate \textbf{T}ransfer) is an \uline{architectural style}, not a wire protocol by itself.
\begin{choices}
  \CorrectChoice True
  \choice False
\end{choices}
\begin{solution}
  REST describes constraints and patterns used \textbf{on top of} protocols such as HTTP; it is not a replacement for HTTP in the way a separate transport protocol would be.
\end{solution}

\question[4]
What is REST?
\begin{choices}
  \choice a application framework
  \CorrectChoice architectural style (convention)
  \choice a programming protocol
  \choice a testing framework
\end{choices}
\begin{solution}
  REST is an \textbf{architectural style} (a set of conventions), not a framework or testing tool.
\end{solution}

\question[4]
\textbf{True or false:} REST technology is preferred in the field because it uses less bandwidth, and is simple and flexible, making it more suitable for internet usage.
\begin{choices}
  \CorrectChoice True
  \choice False
\end{choices}
\begin{solution}
  Text-based formats, statelessness, and standard verbs help REST-style APIs stay practical for web and mobile clients.
\end{solution}

\question[4]
In API design vocabulary, \blank[7em] means giving clients a controlled way to access and manipulate resources (objects) through operations and representations.
\begin{choices}
  \choice Revealing
  \choice Providing
  \choice RESTing
  \CorrectChoice Exposing
\end{choices}
\begin{solution}
  \textbf{Exposing} an API means publishing an interface so clients can access and manipulate resources.
\end{solution}




\question[4]
What is an endpoint?
\begin{solutionorbox}[1.00in]
  An \textbf{endpoint} is a concrete \textbf{HTTP method with path} (and usual headers/body rules) that clients call for one operation; the server answers with a status line, headers, and often a JSON body.
  \begin{lstlisting}[style=http]
GET /api/orders/42 HTTP/1.1
POST /api/orders HTTP/1.1
GET /api/orders?status=open HTTP/1.1
  \end{lstlisting}
  Changing the method or the path (including query) defines a different endpoint even on the same server.
\end{solutionorbox}


\question[4]
What are some best practices for naming resource URIs (paths and queries) using REST?
\begin{solutionorbox}[1.00in]
  Paths should identify \textbf{resources} (nouns). \textbf{What to do} with a resource is expressed with HTTP methods, status codes, and bodies---not by turning the URI path into an RPC-style ``command.'' (In documentation people often say ``URL'' for the full \texttt{https://\ldots} string; the path/query portion is still URI design.)

  \noindent\textbf{Common practices:}
  \begin{itemize}
    \item \textbf{Collections vs.\ items:} plural collection segments (\texttt{/employees}) and stable item paths (\texttt{/employees/42}) are easy to read and to evolve.
    \item \textbf{Nesting:} reflect ownership or containment (\texttt{/departments/3/employees}) when it matches the domain; avoid needlessly deep trees.
    \item \textbf{Methods, not path verbs:} prefer \texttt{POST /orders} over \texttt{/createOrder}; use queries for filters (\texttt{/orders?status=open}) instead of inventing new ``action'' paths.
    \item \textbf{Consistent style:} one casing and separator rule (often lowercase with hyphens) across the whole surface.
    \item \textbf{Versioning:} when breaking changes ship, version deliberately (path, header, or media type) and write the policy down.
  \end{itemize}

  \noindent\textbf{Note:} Clarity and consistency for clients beat strict REST purity in most product APIs.
\end{solutionorbox}

\question[4]
Following REST-style conventions, what is the client most likely trying to do with\\
\texttt{GET http://localhost:8080/user}?
\begin{choices}
  \choice the client is expecting to partially update a user on the server
  \choice the client is expecting to fully update a user on the server
  \CorrectChoice the client expecting to retrieve all users from the server
  \choice the client is expecting to delete a user on the server
\end{choices}
\begin{solution}
  \texttt{GET} is for \textbf{safe reads}. The precise payload (one user vs.\ a list) is defined by that API; the key idea is \textbf{no server-side state change} from a successful \texttt{GET}.
\end{solution}


\question[4]
Following REST-style conventions, what is the client most likely trying to do with\\
\texttt{PUT http://localhost:8080/user}?
\begin{choices}
  \choice The client is expecting to partially update a user on the server
  \CorrectChoice The client is expecting to fully update a user on the server
  \choice The client expecting to retrieve all users from the server
  \choice The client is expecting to delete a user on the server
\end{choices}
\begin{solution}
  \texttt{PUT} is usually \textbf{idempotent full replacement} (or creation at a known URI if the API documents upsert). Compare with \texttt{PATCH} for partial updates.
\end{solution}

\question[4]
Following REST-style conventions, what is the client most likely trying to do with\\
\texttt{PATCH http://localhost:8080/user}?
\begin{choices}
  \CorrectChoice The client is expecting to partially update a user on the server
  \choice The client is expecting to fully update a user on the server
  \choice The client expecting to retrieve all users from the server
  \choice The client is expecting to delete a user on the server
\end{choices}
\begin{solution}
  \texttt{PATCH} carries a \textbf{partial} change set (merge patch, JSON Patch, or vendor-specific rules). Full replacement is more often associated with \texttt{PUT}.
\end{solution}

\question[4]
What is the purpose of the \texttt{PUT} method in RESTful APIs?
\begin{choices}
  \choice Delete data on the server
  \choice Create new data on the server
  \CorrectChoice Replace or update the resource at the target URI
  \choice Retrieve data from the server
\end{choices}
\begin{solution}
  \texttt{PUT} sends a representation to a \textbf{known resource URI}; the server stores it as the resource state---usually a \textbf{full replace}. Some APIs document ``upsert'' (\texttt{PUT} creates if missing); always follow the service's contract.

  \noindent\textbf{Contrast (typical teaching):}
  \begin{itemize}
    \item \textbf{vs.\ \texttt{GET}:} safe read; no body meant to change server state.
    \item \textbf{vs.\ \texttt{POST}:} often creates with server-assigned id or triggers processing on a collection URI, not necessarily idempotent.
    \item \textbf{vs.\ \texttt{PATCH}:} partial update; \texttt{PUT} usually carries the whole resource document.
  \end{itemize}

  \noindent\textbf{Note:} Frameworks vary; correctness is defined by your API documentation and HTTP semantics you advertise.
\end{solution}

\newpage

\question[4]
What is the primary purpose of the \texttt{GET} method in RESTful APIs?
\begin{choices}
  \CorrectChoice Retrieve (read) a representation of a resource
  \choice Create a new resource on the server
  \choice Replace the entire resource at the URI
  \choice Delete the resource at the URI
\end{choices}
\begin{solution}
  \texttt{GET} is a \textbf{safe, read-only} request: fetch a representation (JSON, HTML, etc.). It should not be used to change persistent state; use \texttt{POST}, \texttt{PUT}/\texttt{PATCH}, or \texttt{DELETE} for writes.
\end{solution}

\question[4]
What is a typical purpose of \texttt{POST} to a collection URI (e.g.\ \texttt{POST /orders}) in RESTful APIs?
\begin{choices}
  \choice Delete every order in the collection
  \choice Replace all orders with the request body
  \CorrectChoice Create a new resource (or start processing)---the server often assigns the id and may return \texttt{201 Created}
  \choice List all orders without side effects
\end{choices}
\begin{solution}
  \texttt{POST} is the usual ``\textbf{submit to collection}'' verb: create a subordinate resource, enqueue work, or accept a body whose meaning is defined by the API. It is \textbf{not} assumed idempotent---repeating the same \texttt{POST} may create duplicates unless the API is explicitly designed otherwise.
\end{solution}

\question[4]
What is the purpose of the \texttt{DELETE} method on a resource URI in RESTful APIs?
\begin{choices}
  \choice Retrieve the resource one last time
  \choice Update only selected fields
  \CorrectChoice Remove the resource (or mark it deleted, per API policy)
  \choice Create a copy of the resource
\end{choices}
\begin{solution}
  \texttt{DELETE} asks the server to drop the resource at that URI (or to perform an application-level delete). A \texttt{204 No Content} or \texttt{200 OK} with a body are common success patterns; \texttt{404} applies if it is already gone.
\end{solution}

\question[4]
A client must change \emph{only a few fields} of an existing resource. Which verb is \emph{most} associated with a \textbf{partial update} in REST-style APIs?
\begin{choices}
  \choice \texttt{GET}
  \choice \texttt{POST}
  \CorrectChoice \texttt{PATCH}
  \choice \texttt{PUT} (always, in every API)
\end{choices}
\begin{solution}
  \texttt{PATCH} carries a \textbf{partial} change (JSON Patch, merge patch, or API-specific diff). \texttt{PUT} is often taught as ``replace the whole representation at this URI''; some stacks overload it, but \texttt{PATCH} is the usual label for field-level edits.
\end{solution}


\question[4]
Which HTTP verbs are commonly used in a REST-based architecture?
\begin{choices}
  \choice \texttt{POST}, \texttt{TRACE}, \texttt{PUT}, \texttt{READ}, \texttt{DELETE}
  \CorrectChoice \texttt{PUT}, \texttt{POST}, \texttt{PATCH}, \texttt{GET}, \texttt{DELETE}
  \choice \texttt{PUT}, \texttt{POST}, \texttt{READ}, \texttt{PATCH}, \texttt{TRACE}
  \choice \texttt{PUT}, \texttt{PATCH}, \texttt{POST}, \texttt{UPDATE}, \texttt{DELETE}
\end{choices}
\begin{solution}
  The usual REST teaching set includes \texttt{GET}, \texttt{POST}, \texttt{PUT}, \texttt{PATCH}, and \texttt{DELETE}. \texttt{READ} and \texttt{UPDATE} are not standard HTTP method names; \texttt{TRACE} is rarely emphasized for resource APIs.
\end{solution}


\newpage


\uplevel{\par\textbf{Javalin}\par}

\question[4]
Describe Javalin.
\begin{solutionorbox}[1.00in]
  Javalin is a \textbf{lightweight Java/Kotlin} library or web framework for HTTP \textbf{APIs} and small services: embedded server, little ceremony.

  \begin{lstlisting}[style=java]
  Javalin app = Javalin.create();                    // embedded app / server setup
  app.get("/hello", ctx -> ctx.result("Hi"));       // verb + path + handler; ctx writes body
  app.get("/users/{id}", ctx -> {                   // path params in {braces}
    String id = ctx.pathParam("id");               // read segment from URL
    ctx.json(id);                                  // send JSON (e.g. quoted string here)
  });
  app.start(7070);                                  // listen on port
  \end{lstlisting}
  Filters, WebSockets, and static assets are optional add-ons on the same model. Also, larger ecosystems (e.g.\ Spring) add more integration out of the box. 

  \noindent\textbf{Example HTTP exchanges:}\par
  \noindent The first \texttt{app.get} registers \texttt{GET /hello}. The handler uses \texttt{ctx.result} for a \textbf{plain text} body.\par
  \begin{httpexchange}
\begin{lstlisting}[style=http]
GET /hello HTTP/1.1
Host: localhost:7070
Accept: */*

\end{lstlisting}
\httpmidarrows
\begin{lstlisting}[style=http]
HTTP/1.1 200 OK
Content-Type: text/plain

Hi
\end{lstlisting}
  \end{httpexchange}

  \noindent The second route matches \texttt{GET /users/\{id\}}. \texttt{pathParam} reads the segment; \texttt{ctx.json} sets a \texttt{JSON} body (here the Java \texttt{String} becomes a JSON string literal).\par
  \begin{httpexchange}
\begin{lstlisting}[style=http]
GET /users/42 HTTP/1.1
Host: localhost:7070
Accept: application/json

\end{lstlisting}
\httpmidarrows
\begin{lstlisting}[style=http]
HTTP/1.1 200 OK
Content-Type: application/json

"42"
\end{lstlisting}
  \end{httpexchange}

  \noindent\textbf{Wrong verb:}\par
  \noindent Only \texttt{GET} is registered for these paths. \texttt{POST /hello} does not match any handler and typically yields \texttt{404 Not Found} (unless you add another route or a catch-all).\par
  \begin{httpexchange}
\begin{lstlisting}[style=http]
POST /hello HTTP/1.1
Host: localhost:7070
Content-Type: text/plain
Content-Length: 0


\end{lstlisting}
\httpmidarrows
\begin{lstlisting}[style=http]
HTTP/1.1 404 Not Found
Content-Type: text/plain

Not found
\end{lstlisting}
  \end{httpexchange}

  \noindent\textbf{Note:}
  \texttt{*/*} means ``any media type'' in the request header.
\end{solutionorbox}


\newpage

\question[4]
What is a route handler?
\begin{solutionorbox}[1.00in]
  A \textbf{route handler} is the code the framework runs when an incoming request \textbf{matches} a registered route (HTTP method + path). It reads the request, calls domain logic, and finishes by setting the HTTP result.

  \begin{lstlisting}[style=java]
  app.post("/orders", ctx -> {                // this lambda is the handler for POST /orders
    // read body, validate, call services...
    ctx.status(201).result("created");         // finish: status line + body (often .json(...))
  });
  \end{lstlisting}

  \noindent\textbf{Example HTTP exchange:}\par
  \noindent Client \texttt{POST}s to \texttt{/orders} with a JSON body (example payload below).\par
  \noindent The handler sets status \texttt{201~Created}; the body text comes from \texttt{ctx.result}.\par
  \begin{httpexchange}
\begin{lstlisting}[style=http]
POST /orders HTTP/1.1
Host: localhost:7070
Content-Type: application/json

{"item":"tea"}
\end{lstlisting}
\httpmidarrows
\begin{lstlisting}[style=http]
HTTP/1.1 201 Created
Content-Type: text/plain

created
\end{lstlisting}
  \end{httpexchange}

  \noindent\textbf{Note:} Names vary (\texttt{@GetMapping} method, ``controller action''); the idea is the same: \textbf{one entry point per route pattern}.
\end{solutionorbox}


\newpage


\question[4]
Which method is used to access path parameters in a Javalin handler?
\begin{choices}
  \CorrectChoice \texttt{ctx.pathParam()}
  \choice \texttt{ctx.queryParam()}
  \choice \texttt{ctx.param()}
  \choice \texttt{ctx.getParam()}
\end{choices}
\begin{solution}
  Segments declared in the route template with \texttt{\{name\}} are \textbf{path parameters}. In Javalin you read them with \texttt{ctx.pathParam("name")}.

  \begin{lstlisting}[style=java] 
  app.get("/users/{id}", ctx -> {
    String id = ctx.pathParam("id");   // binds to /users/42 -> "42"
    ctx.result("user " + id);
  });
  \end{lstlisting}

  \noindent\textbf{Example HTTP exchanges:}\par
  \noindent The \texttt{\{id\}} segment in the route pattern is filled by whatever appears in that position in the URL; \texttt{pathParam("id")} returns it as a string.\par
  \begin{httpexchange}
\begin{lstlisting}[style=http]
GET /users/42 HTTP/1.1
Host: localhost:7070
Accept: */*

\end{lstlisting}
\httpmidarrows
\begin{lstlisting}[style=http]
HTTP/1.1 200 OK
Content-Type: text/plain

user 42
\end{lstlisting}
  \end{httpexchange}

  \noindent Same route, different path segment---the handler runs again with a new \texttt{id} value.\par
  \begin{httpexchange}
\begin{lstlisting}[style=http]
GET /users/7 HTTP/1.1
Host: localhost:7070
Accept: */*

\end{lstlisting}
\httpmidarrows
\begin{lstlisting}[style=http]
HTTP/1.1 200 OK
Content-Type: text/plain

user 7
\end{lstlisting}
  \end{httpexchange}

  \noindent\textbf{Wrong verb or no matching route:}\par
  \noindent Only \texttt{GET /users/\{id\}} is registered. \texttt{POST /users/42} (or another verb) does not invoke this handler; neither does \texttt{GET /users} if your app never registered a collection route for that exact path.\par
  \noindent Unmatched method\,+\,path pairs fall through to Javalin's default error handling (\texttt{404 Not Found} is typical unless you add routes or a custom handler).\par
  \begin{httpexchange}
\begin{lstlisting}[style=http]
POST /users/42 HTTP/1.1
Host: localhost:7070
Content-Type: application/json
Content-Length: 2

{}
\end{lstlisting}
\httpmidarrows
\begin{lstlisting}[style=http]
HTTP/1.1 404 Not Found
Content-Type: text/plain

Not found
\end{lstlisting}
  \end{httpexchange}
\end{solution}

\newpage

\question[4]
Which \texttt{Context} method reads values from the \textbf{query string} (the part after \texttt{?})? \\
\textit{e.g.} \texttt{/search?q=tea\&sort=name}
\begin{choices}
  \choice \texttt{ctx.pathParam()}
  \CorrectChoice \texttt{ctx.queryParam()}
  \choice \texttt{ctx.param()} 
  \choice \texttt{ctx.getParam()}
\end{choices}
\begin{solution}
  Query parameters are \texttt{name=value} pairs in the URL after \texttt{?}. Javalin exposes them with \texttt{ctx.queryParam("name")} (optional default overload when the key is absent).

  \begin{lstlisting}[style=java]
  app.get("/search", ctx -> {
    String q = ctx.queryParam("q");           // /search?q=tea -> "tea"
    String sort = ctx.queryParam("sort", "id"); // default if missing
    ctx.result(q + " sorted by " + sort);
  });
  \end{lstlisting}

  \noindent\textbf{Example HTTP exchanges:}\par
  \noindent Both \texttt{q} and \texttt{sort} present---\texttt{ctx.queryParam} reads each value from the query string.\par
  \begin{httpexchange}
\begin{lstlisting}[style=http]
GET /search?q=tea&sort=name HTTP/1.1
Host: localhost:7070
Accept: */*

\end{lstlisting}
\httpmidarrows
\begin{lstlisting}[style=http]
HTTP/1.1 200 OK
Content-Type: text/plain

tea sorted by name
\end{lstlisting}
  \end{httpexchange}

  \noindent Missing \texttt{sort}: the two-argument call \texttt{queryParam("sort", "id")} supplies the default \texttt{id}.\par
  \begin{httpexchange}
\begin{lstlisting}[style=http]
GET /search?q=tea HTTP/1.1
Host: localhost:7070
Accept: */*

\end{lstlisting}
\httpmidarrows
\begin{lstlisting}[style=http]
HTTP/1.1 200 OK
Content-Type: text/plain

tea sorted by id
\end{lstlisting}
  \end{httpexchange}

  \noindent\textbf{Wrong verb or no matching route:}\par
  \noindent Only \texttt{GET /search} is registered above. A \texttt{POST} (or \texttt{PUT}, \texttt{DELETE}, \ldots) to the same path, or \texttt{GET} to a path you never registered, does \emph{not} run that lambda.\par
  \noindent Javalin's default is an error response for unmatched method\,+\,path pairs; you usually see \texttt{404 Not Found} unless you add another route or a custom error handler.\par
  \begin{httpexchange}
\begin{lstlisting}[style=http]
POST /search HTTP/1.1
Host: localhost:7070
Content-Type: application/json
Content-Length: 2

{}
\end{lstlisting}
\httpmidarrows
\begin{lstlisting}[style=http]
HTTP/1.1 404 Not Found
Content-Type: text/plain

Not found
\end{lstlisting}
  \end{httpexchange}
\end{solution}

\newpage

\question[4]
In Javalin, what is the purpose of the context object parameter in route handlers?
\begin{choices}
  \choice It represents the request body sent by the client
  \CorrectChoice It provides access to the HTTP request and sets the response objects
  \choice It defines the route path for the handler
  \choice It specifies the HTTP method for the handler
\end{choices}
\begin{solution}
  \texttt{Context} (\texttt{ctx}) is created \textbf{per request}: one object that exposes the \textbf{incoming} HTTP message (path, query, headers, body, cookies, \ldots) and helpers to build the \textbf{outgoing} response (status, headers, body, redirect).

  \begin{lstlisting}[style=java]
  app.get("/demo", ctx -> {
    String q = ctx.queryParam("q");           // read request inputs
    ctx.header("X-App", "demo");              // response header
    ctx.status(200).result("search=" + q);    // status + body
  });
  \end{lstlisting}

  \noindent\textbf{Example HTTP exchange:}\par
  \noindent The client sends \texttt{q} in the query string; \texttt{ctx} supplies that value to the handler. The handler adds a response header, sets \texttt{200}, and writes the body built from \texttt{q}.\par
  \begin{httpexchange}
\begin{lstlisting}[style=http]
GET /demo?q=tea HTTP/1.1
Host: localhost:7070
Accept: */*

\end{lstlisting}
\httpmidarrows
\begin{lstlisting}[style=http]
HTTP/1.1 200 OK
Content-Type: text/plain
X-App: demo

search=tea
\end{lstlisting}
  \end{httpexchange}
\end{solution}

\newpage

\question[4]
In Javalin, how do you read the raw \textbf{request body} (for example JSON as text) inside a handler?
\begin{choices}
  \choice The \texttt{ctx} parameter is the request body
  \CorrectChoice \texttt{ctx.body()} (often followed by parsing, e.g.\ \texttt{bodyAsClass(...)} for JSON)
  \choice \texttt{ctx.pathParam(...)}
  \choice \texttt{ctx.queryParam(...)}
\end{choices}
\begin{solution}
  The body is a \textbf{stream of bytes} the client sent; \texttt{ctx} is the whole request/response facade. Use \texttt{ctx.body()} for the payload as a \texttt{String}, or mapping helpers for typed JSON.

  \begin{lstlisting}[style=java]
  app.post("/items", ctx -> {
    String raw = ctx.body();                  // raw body text
    // ctx.bodyAsClass(Item.class);           // typed JSON (method name varies by Javalin version)
    ctx.status(raw.isEmpty() ? 400 : 201).result("ok");
  });
  \end{lstlisting}

  \noindent\textbf{Example HTTP exchanges:}\par
  \noindent \texttt{ctx.body()} returns the request payload as text (here JSON). The handler treats an empty body as a client error; otherwise it answers \texttt{201~Created}.\par
  \begin{httpexchange}
\begin{lstlisting}[style=http]
POST /items HTTP/1.1
Host: localhost:7070
Content-Type: application/json

{"name":"mug","qty":2}
\end{lstlisting}
\httpmidarrows
\begin{lstlisting}[style=http]
HTTP/1.1 201 Created
Content-Type: text/plain

ok
\end{lstlisting}
  \end{httpexchange}

  \noindent Same route with \textbf{no} message body: \texttt{ctx.body()} is the empty string, so the sample handler picks \texttt{400~Bad Request}.\par
  \begin{httpexchange}
\begin{lstlisting}[style=http]
POST /items HTTP/1.1
Host: localhost:7070
Content-Length: 0


\end{lstlisting}
\httpmidarrows
\begin{lstlisting}[style=http]
HTTP/1.1 400 Bad Request
Content-Type: text/plain

ok
\end{lstlisting}
  \end{httpexchange}
\end{solution}

\question[4]
Where is the \textbf{route path} (e.g.\ \texttt{/orders/\{id\}}) specified for a Javalin handler?
\begin{choices}
  \CorrectChoice In the route registration call---first argument to \texttt{app.get(...)}, \texttt{app.post(...)}, etc.
  \choice Inside the handler by assigning to \texttt{ctx}
  \choice Only in a separate XML configuration file
  \choice Implicitly from the name of the handler method
\end{choices}
\begin{solution}
  Javalin wires \texttt{method + path pattern + handler} together at registration time. The string \texttt{"/orders/\{id\}"} belongs next to \texttt{app.get} or \texttt{app.post}, not inside \texttt{ctx}.

  \begin{lstlisting}[style=java]
  app.get("/orders/{id}", ctx -> {             // path pattern fixed here
    ctx.result(ctx.pathParam("id"));          // ctx only reads matched values
  });
  \end{lstlisting}
\end{solution}


\newpage

\question[4]
What determines which \textbf{HTTP method} (e.g.\ \texttt{GET} vs.\ \texttt{POST}) a given handler accepts?
\begin{choices}
  \choice The handler sets it with \texttt{ctx.method(...)} before returning
  \CorrectChoice Which registrar you use: \texttt{app.get}, \texttt{app.post}, \texttt{app.put}, \texttt{app.delete}, \ldots
  \choice The filename of the Java source file
  \choice The return type of the lambda
\end{choices}
\begin{solution}
  Each fluent call registers one verb for that path pattern. A \texttt{GET} request never runs a handler attached only with \texttt{app.post}; you need a matching registration (or a shared handler pattern you configure per verb).

  \begin{lstlisting}[style=java]
  app.get("/ping", ctx -> ctx.result("GET"));   // handles GET /ping
  app.post("/ping", ctx -> ctx.result("POST")); // same path, different method -> different handler
  \end{lstlisting}
\end{solution}

\question[4]
Which of the following is a valid way to define a Javalin route handler for handling HTTP POST requests?
\begin{choices}
  \choice \texttt{@PostMapping("/path", ctx -> \{ \})}
  \choice \texttt{@Route("/path", ctx -> \{ \})}
  \CorrectChoice \texttt{post("/path", ctx -> \{ \});}
  \choice \texttt{handler.post("/path", ctx -> \{ \})}
\end{choices}
\begin{solution}
  Javalin registers routes with \textbf{fluent calls} on the \texttt{Javalin} instance. 
  The pattern is \texttt{app.post(path, ctx -> \{ \ldots\ \})}: HTTP verb as the method name, 
  path string, then a lambda that receives the per-request \texttt{Context}. 

  \texttt{@PostMapping} is \textbf{Spring MVC}: annotations on controller methods, not Javalin's registration style. \texttt{@Route(...)} with a lambda is not how Javalin wires handlers (other stacks use different \texttt{@Route} meanings). \texttt{handler.post(...)} suggests the wrong object; you call \texttt{post} on the \texttt{Javalin} app, not on an unrelated \texttt{handler} variable.

  \begin{lstlisting}[style=java]
  app.post("/path", ctx -> {
    ctx.status(201).result("ok");   // POST /path -> this lambda
  });
  \end{lstlisting}
\end{solution}

\newpage

\question[4]
How do you send a JSON response with status code 201 (Created) in Javalin?
\begin{choices}
  \CorrectChoice \texttt{ctx.json(data).status(201)}
  \choice \texttt{ctx.result(data).status(201)}
  \choice \texttt{ctx.send(data, "application/json", 201)}
  \choice \texttt{ctx.returnJson(data).status(201)}
\end{choices}
\begin{solution}
  Use \texttt{ctx.json(data)} so Javalin \textbf{serializes} \texttt{data} to JSON (via the configured mapper, often Jackson) and sets an appropriate \texttt{Content-Type} (typically \texttt{application/json}). Chain \texttt{.status(201)} to send \textbf{201~Created}, the usual success code after a \texttt{POST} that creates a resource. The fluent methods return \texttt{ctx}, so \texttt{ctx.json(data).status(201)} and \texttt{ctx.status(201).json(data)} are both common orderings.

  \texttt{ctx.result(\ldots)} writes the argument as the \textbf{raw response body} (often \texttt{text/plain}); it does not treat a Java object as JSON unless you stringify it yourself, so it is the wrong tool for a typed JSON payload. \texttt{ctx.send(\ldots)} and \texttt{returnJson} are not the usual \texttt{Context} API in Javalin (method names and signatures differ from these distractors).

  \begin{lstlisting}[style=java]
  app.post("/items", ctx -> {
    Item created = service.create(/* ... parse body ... */);
    ctx.json(created).status(201);   // JSON body + 201 Created
  });
  \end{lstlisting}

  \noindent\textbf{Example HTTP exchange:}\par
  \noindent Client \texttt{POST}s JSON; the handler builds \texttt{created}, then \texttt{ctx.json} becomes the response body and \texttt{status(201)} sets the status line.\par
  \begin{httpexchange}
\begin{lstlisting}[style=http]
POST /items HTTP/1.1
Host: localhost:7070
Content-Type: application/json

{"name":"mug","qty":1}
\end{lstlisting}
\httpmidarrows
\begin{lstlisting}[style=http]
HTTP/1.1 201 Created
Content-Type: application/json

{"id":7,"name":"mug","qty":1}
\end{lstlisting}
  \end{httpexchange}
\end{solution}

\question[4]
What is the purpose of Javalin's \texttt{ExceptionHandler} interface?
\begin{choices}
  \choice It provides methods for generating custom exceptions.
  \CorrectChoice It defines a contract for classes that handle exceptions in Javalin applications.
  \choice It automatically handles all exceptions thrown by route handlers.
  \choice It enables logging of exceptions without affecting application flow.
\end{choices}
\begin{solution}
  \texttt{ExceptionHandler} is a \textbf{functional interface}: one method that takes the thrown exception and the request \texttt{Context}. You implement it (often with a lambda) to turn failures into \textbf{HTTP status lines, JSON error bodies, headers}, or side effects such as logging---instead of leaking a raw stack trace to the client.

  Register handlers on the \texttt{Javalin} instance with \texttt{app.exception(ExceptionType.class, handler)}. Javalin picks the \textbf{most specific} matching type for an uncaught exception from a route or filter. Anything you do not register still falls through to the framework's \textbf{default} error behavior (commonly \texttt{500} with a generic body unless you override it with a broad \texttt{Exception.class} handler).

  \begin{lstlisting}[style=java]
  app.exception(IllegalArgumentException.class, (e, ctx) -> {
    ctx.status(400).result("bad request: " + e.getMessage());
  });
  app.exception(Exception.class, (e, ctx) -> {   // broader fallback
    ctx.status(500).result("internal error");
  });
  \end{lstlisting}
\end{solution}

\newpage

\uplevel{\par\textbf{HTTP and REST (practice bank)}\par}

\question[4]
\blank[3.5em] retrieves a representation of a resource.
\begin{choices}
  \CorrectChoice \texttt{GET}
  \choice \texttt{POST}
  \choice \texttt{DELETE}
  \choice \texttt{PUT}
\end{choices}
\begin{solution}
  \textbf{GET} requests a read-only representation (safe/idempotent in typical REST teaching). \texttt{POST} often creates or submits; \texttt{PUT}/\texttt{PATCH} replace or update; \texttt{DELETE} removes.
\end{solution}

\question[4]
What does this HTTP response mean that was returned from a server?
\begin{lstlisting}[style=http]
HTTP/1.1 404 Not Found
Date: Sun, 18 Oct 2012 10:36:20 GMT
Server: Apache/2.2.14 (Win32)
Content-Length: 230
Connection: Closed
Content-Type: text/html; charset=iso-8859-1


    404 Not Found

    Not Found
    The requested URL /t.html was not found on this server.
\end{lstlisting}
\begin{choices}
  \CorrectChoice There was a client-side error
  \choice There was a server-side error
  \choice The response returned the results we expected
  \choice There was a redirect to another web page
\end{choices}
\begin{solution}
  \texttt{404 Not Found} is a \textbf{4xx} client error: the server understood the request but the resource path was not found (from the server's point of view).
\end{solution}

\newpage


\question[4]
What, if anything, is wrong with this endpoint when designing a REST API?
\begin{center}
  \texttt{mylibrarydemo.com/api/getAllBooks}
\end{center}
\begin{choices}
  \CorrectChoice RESTful endpoints should refer to nouns, not verbs. The HTTP method in the request should instead specify the action to be taken.
  \choice There is nothing wrong with this RESTful endpoint
  \choice Query or path parameters should be used to retrieve all books
  \choice Retrieval of a resource (a \texttt{GET} operation) is not allowed in a REST API
\end{choices}
\begin{solution}
  The issue is the verb in the path (\texttt{getAllBooks}). In REST, the \textbf{URI names the resource} (noun), and the \textbf{HTTP method expresses the action}.

  \textbf{Less RESTful:}
  \begin{itemize}
    \item \texttt{GET /api/getAllBooks}
    \item \texttt{POST /api/createBook}
    \item \texttt{DELETE /api/deleteBook/42}
  \end{itemize}

  \textbf{REST-style equivalent:}
  \begin{itemize}
    \item \texttt{GET /api/books} \quad (list all books)
    \item \texttt{GET /api/books/42} \quad (fetch one book)
    \item \texttt{POST /api/books} \quad (create a book)
    \item \texttt{PUT /api/books/42} or \texttt{PATCH /api/books/42} \quad (update)
    \item \texttt{DELETE /api/books/42} \quad (delete)
  \end{itemize}

  So the precise fix is: use \texttt{/api/books} (and \texttt{/api/books/\{id\}}) as noun-based endpoints, and rely on HTTP verbs for behavior.
\end{solution}

\question[4]
This is an HTTP response received from a server. What does this mean?
\begin{lstlisting}[style=http]
HTTP/1.1 200 OK
Access-Control-Allow-Origin: *
Content-Type: text/html

    Hello World
\end{lstlisting}
\begin{choices}
  \choice There was an error on the server side
  \choice The server redirected the client to another URL
  \choice There was an error with the request itself and it could not be completed
  \CorrectChoice The request was completed successfully and the server returned an HTML page to the client
\end{choices}
\begin{solution}
  \texttt{200 OK} means success. \texttt{Content-Type: text/html} and the body indicate an HTML response. (Some banks incorrectly also mark a ``request error'' choice; that contradicts \texttt{200}.)
\end{solution}


\newpage

\question[4]
You want a REST API so clients can create a new book by sending JSON like this:
\begin{lstlisting}[style=java]
{
    "ISBN": 182037184371,
    "Title": "Adventures of Huckleberry Finn",
    "Author": "Mark Twain"
}
\end{lstlisting}
What kind of HTTP method should be used, and where does this data belong?
\begin{choices}
  \choice Use a \texttt{GET} request; data should be in the body
  \choice Use a \texttt{POST} request; data should be in the header
  \CorrectChoice Use a \texttt{POST} request; data should be in the body
  \choice Use a \texttt{GET} request; data should be in the header
\end{choices}
\begin{solution}
  Creating a resource usually uses \textbf{POST} (sometimes \texttt{PUT} to a known URL) with the representation in the \textbf{message body} (often \texttt{application/json}). Headers carry metadata, not the full JSON document.
\end{solution}

\endgradingrange{csarest}
